\scanpage{78}
\mspage{27}

\uline{Dém}. Pour expliciter la condition a), il faut \deleted{\unsure{o}} \deleted{expliciter les fibres de \unsure{ce mor}} expliciter la catégorie $  {}_{\unsure{\xi|\zeta}_0}\backslash(X_y)$, pour un objet quelconque \inserted{$\xi_0=$} $(x_0,u_0)$ de $X_{/y}$, \deleted{i.e. que une flèche} \deleted{quelconque de but $y$ de $X$} \inserted{où} $x_0\in\mathrm{Ob}\,X$, \deleted{\unsure{et}} \[u_0:y_0\longrightarrow y_1=y\quad\text{de }Y,\text{ avec }y_0=f(x_0).\] Cette catégorie $  {}_{\xi_0}\backslash X_{y_1}$ s'explicite alors comme la catégorie des $x_1$ dans $X_{y_1}$, \uline{munis} d'une flèche \inserted{de $\xi_0$} dans \deleted{\unsure{la}} l'image \inserted{$\xi_1=(x_1,\mathrm{id}_{f(x_1)}\unsure{=|\Rightarrow}\mathrm{id}_{y_1})$} de $x_1$ dans $X_{/y_1}$ de $\xi_0$, i.e. d'une flèche $u$ de $x_0$ dans $x_1$, \diagram{\begin{tikzcd}[ampersand replacement=\&,column sep=small,row sep=small] x_1\arrow[d,dash]\&\&x_0\arrow[ll,"u"']\\ y_1\&y_1\arrow[l,dash,"{\mathrm{id}}"']\&\\ \&\&y_0\arrow[uu]\arrow[ul,"{u_0}"']\arrow[ull,"{f(u)}"] \end{tikzcd}} telle que $f(u)$ donne lieu au triangle comm. dans $Y$, i.e. $f(u)=u_0$. Donc c'est la catégorie notée $X(x_0,u_0)$ dans d). On voit tout de suite que \deleted{c'est} \inserted{elle est} aussi \inserted{iso à} la fibre de $  {}_{x_0}\backslash X\mathrel{\relbar}{}_{y_0}\backslash Y$ en l'élément $u_0$ de $ {y_0}\backslash Y$, dont il est question dans b). D'où l'équivalence de a) b) \unsure{c)|d)}.

D'autre part, la donnée de $u_0$ dans $Y$ équivaut à celle d'un foncteur $\mathord{\tikz[baseline=0pt,x=1ex,y=1ex]{\draw[line width=.55pt,line join=round] (0,0)--(.78,1.4)--(1.53,0)--cycle; \draw[line width=.55pt] (.25,0)--(.90,1.18);}}^1\mathrel{\relbar}Y$, \unsure{comme} dans c), et pour exprimer alors que $\underset{\substack{\mathord{\tikz[baseline=0pt,x=1ex,y=1ex]{\draw[line width=.45pt] (0,0) .. controls (-.28,.45) and (.28,.65) .. (0,1); \draw[line width=.45pt] (.30,0)--(.30,1);}}\\X_{y_1}}}{X'_1}\longrightarrow X'$ est coasphérique, il suffit de \deleted{\unsure{vérifier}} \inserted{l'exprimer} que les $  {}_{x'}\backslash X'_1$ \inserted{($x'\in\mathrm{Ob}\,X'$)} sont $W$-asphérique. Or c'est trivial si $x'\in\mathrm{Ob}\,X'_1$ (objet initial), donc il suffit de le vérifier pour $x'$ \deleted{$\in$} \inserted{dans} $X'\smallsetminus X'_1=X'_0\simeq X_{y_0}$. \deleted{\unsure{Et}} \inserted{Alors $x'$ s'identifie à un élément $x_0$ de $X_{y_0}$,} on vérifie tout de suite que $  {}_{x_0}\backslash X'_1$ est isom à la catégorie $X(x_0,u_0)$ de d). \unsure{Céla|Cela} achève la dém.

% archive sheet number: 78
